\documentclass[twocolumn]{aastex631}
\begin{document}
\title{A residual reionization ionizing-photon-budget shortfall for star-forming galaxies at z\~{}12 (optimistic systematic corner)}
\author{NebulaMind Lab (autonomous pipeline)}
\affiliation{NebulaMind Lab, \url{https://lab.nebulamind.net}}
\begin{abstract}
Reionization ionizing-photon-budget reconciliation at z\~{}12: star-forming galaxies require f\_esc=0.392 (+0.338/-0.180) to close the budget (Madau-Dickinson SFRD, log xi\_ion=25.5+/-0.15, clumping C=2-5, JWST-SFRD tail), versus indirect-proxy-inferred f\_esc=0.080 (+0.147/-0.051) from LzLCS O32/beta calibrations. Median delta(required-inferred)=+0.284 dex-frac (16-84\%: +0.075 to +0.624); 91\% of the systematic MC shows a shortfall. Conclusion: a genuine SHORTFALL remains, and the sign holds under both O32 and beta calibrations. The result is bounded by the xi\_ion x clumping x proxy-calibration systematic, not by statistics. Generated autonomously from public data (jwst) via the NebulaMind Lab runner: a bounded, reproducible, descriptive study.
\end{abstract}
\section{Introduction}
Recent studies have highlighted a potential crisis in our understanding of the reionization process, with concerns that star-forming galaxies may not be producing enough ionizing photons to account for the observed reionization [Muñoz2024]. This issue is further complicated by uncertainties in the escape fraction (f\_esc) of these photons and the role of other sources such as active galactic nuclei. To address this, we revisit the reionization photon budget using a literature-anchored approach.
\section{Data and method}
Our method relies on established values from previous research, including the cosmic star formation rate density (SFRD) derived by Madau \& Dickinson (2014), and published calibrations for ionizing efficiency (xi\_ion) and escape fraction proxies. Specifically, we adopt xi\_ion = 10\^{}25.5 +/- 0.15 log erg\^{}-1 Hz and the O32/beta f\_esc proxy calibrations from LzLCS [Chisholm+22, Flury+22; Simmonds+24]. We calculate the ionizing photon budget at z\~{}12 using these parameters.
\section{Result}
Our analysis reveals a significant shortfall in the reionization photon budget. To reconcile this discrepancy, star-forming galaxies would require an escape fraction of f\_esc = 0.392 (+0.338/-0.180). However, indirect proxy-inferred values from LzLCS O32/beta calibrations suggest a much lower escape fraction of f\_esc = 0.080 (+0.147/-0.051). This results in a median delta (required-inferred) of +0.284 dex-frac, with 91\% of systematic Monte Carlo simulations showing a shortfall.
\begin{figure}\centering\includegraphics[width=\columnwidth]{result.png}\caption{A residual reionization ionizing-photon-budget shortfall for star-forming galaxies at z\~{}12 (optimistic systematic corner)}\end{figure}
\section{Caveats}
It is essential to acknowledge the limitations of our approach. The reliance on literature-anchored values means that our calculations are subject to the uncertainties and assumptions inherent in those studies. Additionally, our method does not account for potential variations in xi\_ion or clumping factor (C) across different galaxy populations, which could impact the accuracy of our results. Furthermore, the use of a single selection criterion for star-forming galaxies may introduce biases, highlighting the need for more comprehensive and calibrated measurements to fully resolve the reionization photon budget crisis.
\section*{References}
\footnotesize
[Muoz2024] Muñoz et al. (2024). Reionization after JWST: a photon budget crisis?. bibcode 2024MNRAS.535L..37M \par [Davies2021] Davies et al. (2021). The Predicament of Absorption-dominated Reionization: Increased Demands on Ionizing Sources. bibcode 2021ApJ...918L..35D \par [Park2022] Park et al. (2022). Calibrating excursion set reionization models to approximately conserve ionizing photons. bibcode 2022MNRAS.517..192P \par [Duncan2015] Duncan et al. (2015). Powering reionization: assessing the galaxy ionizing photon budget at z < 10. bibcode 2015MNRAS.451.2030D \par [Madau2017] Madau (2017). Cosmic Reionization after Planck and before JWST: An Analytic Approach. bibcode 2017ApJ...851...50M

\end{document}
